\begin{center}
\begin{tikzpicture}[
  x=1cm,
  y=1cm,
  font=\scriptsize,
  levelband/.style={fill=blue!3, rounded corners=1pt},
  treeedge/.style={draw=black!25, line width=0.35pt},
  treenode/.style={circle, draw=black!42, fill=black!4,
                   minimum size=2.8mm, inner sep=0pt, line width=0.35pt},
  leafnode/.style={treenode, minimum size=1.8mm, fill=orange!9},
  claim/.style={draw=blue!55!black, fill=blue!4, rounded corners=1.5pt,
                minimum width=7.15cm, minimum height=0.42cm,
                inner sep=1.1pt, line width=0.45pt},
  rowsc/.style={draw=blue!55!black, fill=blue!10, rounded corners=1.2pt,
                minimum width=2.20cm, minimum height=0.48cm,
                align=center, inner sep=1pt, line width=0.35pt},
  colsc/.style={draw=orange!70!black, fill=orange!12, rounded corners=1.2pt,
                minimum width=2.45cm, minimum height=0.48cm,
                align=center, inner sep=1pt, line width=0.35pt},
  linefold/.style={draw=black!45, fill=black!4, rounded corners=1.2pt,
                   minimum width=1.55cm, minimum height=0.48cm,
                   align=center, inner sep=1pt, line width=0.35pt},
  claimarr/.style={-{Latex[length=1.5mm]}, draw=blue!50!black,
                   line width=0.45pt},
  tinyarr/.style={-{Latex[length=1.2mm]}, draw=black!45,
                   line width=0.35pt},
  rowdetail/.style={draw=blue!55!black, fill=blue!7, rounded corners=1.5pt,
                    text width=3.80cm, minimum height=1.68cm,
                    align=center, inner sep=3pt, line width=0.4pt},
  coldetail/.style={draw=orange!70!black, fill=orange!8,
                    rounded corners=1.5pt, text width=3.75cm,
                    minimum height=1.68cm, align=center,
                    inner sep=3pt, line width=0.4pt},
  folddetail/.style={draw=black!45, fill=black!3, rounded corners=1.5pt,
                     text width=3.35cm, minimum height=1.68cm,
                     align=center, inner sep=3pt, line width=0.4pt},
  boundary/.style={draw=orange!70!black, fill=orange!7,
                   rounded corners=1.5pt, align=center,
                   inner sep=3pt, line width=0.45pt},
  dcap/.style={font=\tiny, text=black!58, align=center},
]

% One continuous swimlane: every level table is horizontally aligned with the
% verifier claim about that table.  Operations in the gap reduce that claim to
% the claim on the next level.
\node[dcap, text width=4.1cm] at (2.13,0.70)
  {one representative column tree; the whole tree is repeated $2^s$ times};
\node[dcap, text width=7.15cm] at (8.28,0.70)
  {each level claim is followed by its row-prefix sumcheck, column-index sumcheck,
   and line fold, side by side};

\draw[draw=black!22, dashed, rounded corners=2pt, line width=0.35pt]
  (0.10,0.30) rectangle (4.22,-4.62);
\node[dcap, fill=white, inner sep=0.7pt] at (2.16,-4.62)
  {one complete product tree};

\foreach \lev in {0,...,4} {
  \pgfmathtruncatemacro{\numnodes}{2^\lev}
  \pgfmathtruncatemacro{\lastnode}{\numnodes-1}
  \pgfmathsetmacro{\yy}{-1.10*\lev}
  \path[levelband] (0.18,\yy-0.17) rectangle (4.14,\yy+0.17);
  \foreach \idx in {0,...,\lastnode} {
    \pgfmathsetmacro{\xx}{0.24+3.84*(\idx+0.5)/\numnodes}
    \coordinate (gkr-\lev-\idx) at (\xx,\yy);
  }
  \ifnum\lev=4
    \node[dcap, anchor=east] at (0.06,\yy) {$L_4=\Lambda$};
  \else
    \node[dcap, anchor=east] at (0.06,\yy) {$L_{\lev}$};
  \fi
}

% Recursive visual order keeps the binary tree legible.  The surrounding prose
% states that the implementation's flat tables instead use low-bit-first order.
\foreach \lev in {0,...,3} {
  \pgfmathtruncatemacro{\numnodes}{2^\lev}
  \pgfmathtruncatemacro{\lastnode}{\numnodes-1}
  \pgfmathtruncatemacro{\childlev}{\lev+1}
  \foreach \idx in {0,...,\lastnode} {
    \pgfmathtruncatemacro{\leftchild}{2*\idx}
    \pgfmathtruncatemacro{\rightchild}{2*\idx+1}
    \draw[treeedge] (gkr-\lev-\idx) -- (gkr-\childlev-\leftchild);
    \draw[treeedge] (gkr-\lev-\idx) -- (gkr-\childlev-\rightchild);
  }
}

\foreach \lev in {0,...,4} {
  \pgfmathtruncatemacro{\numnodes}{2^\lev}
  \pgfmathtruncatemacro{\lastnode}{\numnodes-1}
  \foreach \idx in {0,...,\lastnode} {
    \ifnum\lev=4
      \node[leafnode] at (gkr-\lev-\idx) {};
    \else
      \node[treenode] at (gkr-\lev-\idx) {};
    \fi
  }
}

\draw[-{Latex[length=1.7mm]}, draw=black!38, line width=0.45pt]
  (-0.36,-4.22) -- (-0.36,-0.18);
\node[dcap, rotate=90, anchor=center] at (-0.68,-2.20)
  {products computed upward};

% Claims live at the same heights as the product-tree levels.
\node[claim] (E0) at (8.28, 0.00)
  {$E_0^{\mathrm{GKR}}:\ \mle{L_0}(\zcol^{(0)})=e_0$};
\node[claim] (E1) at (8.28,-1.10)
  {$E_1^{\mathrm{GKR}}:\ \mle{L_1}(\zrow^{(1)},\zcol^{(1)})=e_1$};
\node[claim] (E2) at (8.28,-2.20)
  {$E_2^{\mathrm{GKR}}:\ \mle{L_2}(\zrow^{(2)},\zcol^{(2)})=e_2$};
\node[claim] (E3) at (8.28,-3.30)
  {$E_3^{\mathrm{GKR}}:\ \mle{L_3}(\zrow^{(3)},\zcol^{(3)})=e_3$};
\node[claim] (E4) at (8.28,-4.40)
  {$E_4^{\mathrm{GKR}}:\ \mle{L_4}(\zrow^{(4)},\zcol^{(4)})=e_4$};

\foreach \k/\next/\rowrounds/\midy in {
  0/1/0/-0.55,
  1/2/1/-1.65,
  2/3/2/-2.75,
  3/4/3/-3.85} {
  \node[rowsc] (R\k) at (5.70,\midy)
    {row-prefix sumcheck\\$\rowrounds$ rounds};
  \node[colsc] (C\k) at (8.28,\midy)
    {column-index sumcheck\\$s$ rounds};
  \node[linefold] (X\k) at (10.82,\midy)
    {$\chi_{\k}$ line fold};
  \draw[claimarr] (E\k.south west) -- (R\k.north west);
  \draw[tinyarr] (R\k) -- (C\k);
  \draw[tinyarr] (C\k) -- (X\k);
  \draw[claimarr] (X\k.south east) -- (E\next.north east);
}

% Light guides make the level-table / claim correspondence explicit without
% implying that a single selected tree path is being verified.
\foreach \lev in {0,...,4} {
  \draw[draw=black!20, densely dotted, line width=0.3pt]
    (4.22,-1.10*\lev) -- (E\lev.west);
}

% Expand one transition into its exact three consecutive relations.  The boxes
% are horizontal because these are stages of one claim reduction, not functions
% named A and B.
\node[rowdetail] (rowdetail) at (2.05,-6.30) {%
  \textbf{Row-prefix sumcheck} \ ($k$ rounds)\\[2pt]
  $\begin{aligned}
    e_k&=\displaystyle\sum_{x\in\bits^k}
      \eq(x,\zrow^{(k)})\\[-1pt]
      &\qquad\cdot G_k(x),\\[2pt]
    e_{k,\mathrm{row}}^{\mathrm{end}}
      &=\eq(r_{\mathrm{row}}^{(k)},\zrow^{(k)})s_k
  \end{aligned}$
};

\node[coldetail] (coldetail) at (6.25,-6.30) {%
  \textbf{Column-index sumcheck} \ ($s$ rounds)\\[2pt]
  $\begin{aligned}
    s_k&=\displaystyle\sum_{c\in\bits^s}
      \eq(c,\zcol^{(k)})\\[-1pt]
      &\qquad\cdot E_k(c)O_k(c),\\[2pt]
    e_{k,\mathrm{col}}^{\mathrm{end}}
      &=\eq(r_{\mathrm{col}}^{(k)},\zcol^{(k)})a_kb_k
  \end{aligned}$
};

\node[folddetail] (folddetail) at (10.34,-6.30) {%
  \textbf{Absorb the child pair; then fold}\\[2pt]
  $\begin{gathered}
    e_{k+1}=a_k+\chi_k(a_k+b_k),\\[2pt]
    \zrow^{(k+1)}=(r_{\mathrm{row}}^{(k)},\chi_k),\\[1pt]
    \zcol^{(k+1)}=r_{\mathrm{col}}^{(k)}
  \end{gathered}$
};

\draw[claimarr] (rowdetail) -- (coldetail);
\draw[claimarr] (coldetail) -- (folddetail);

\node[boundary, text width=12.05cm] at (6.14,-8.62) {%
  $\displaystyle
  \begin{aligned}
    e_4+1
      &=\sum_{i\in\bits^4}\sum_{c\in\bits^s}
        \eq(i,\zrow^{(4)})\eq(c,\zcol^{(4)})
        \bigl(\gen^{\gam[i]}-1\bigr)\pitwo(H[i,c])\\[-1pt]
      &\hspace{5.3cm}\longrightarrow\quad
        \text{row-only pre-sumcheck}.
  \end{aligned}$
};
\end{tikzpicture}
\captionof{figure}{A depth-$4$ merged GKR forest drawn as one claim-reduction
swimlane.  Each level table $L_k$ is horizontally aligned with its open claim
$E_k^{\mathrm{GKR}}$.  Between levels, the row-prefix and column-index
sumchecks are displayed side by side, followed by the line fold that creates
the next claim.  The row-prefix sumcheck has $k$ degree-$3$ rounds at transition
$k$ and is empty at the root; the column-index sumcheck always has $s$ rounds.
The displayed tree uses recursive visual order rather than the implementation's
low-bit-first flat-table order.  Across all four transitions there are $6+4s$
degree-$3$ sumcheck rounds and four line folds.}
\label{fig:gkr-depth4}
\end{center}
